\documentclass{beamer}
\usetheme{Madrid}
\usecolortheme{default}
\usepackage[utf8]{inputenc}
\usepackage{amsmath}
\usepackage{amsfonts}
\usepackage{amssymb}
\usepackage{tikz}

\title{CRISC Domain D - 100 More Impossible Hard Questions}
\subtitle{Information Technology and Security - Expert Level Part 3}
\author{CRISC Mastery Series}
\date{\today}

\begin{document}
	
	\begin{frame}
		\titlepage
	\end{frame}
	
	\begin{frame}{Warning}
		\begin{alertblock}{Extreme Difficulty}
			Domain D is technical but the exam tests judgment, not trivia.\\
			Every option will look plausible.\\
			Think like a risk practitioner advising on technology, not a technician.
		\end{alertblock}
	\end{frame}
	
	\begin{frame}{Outline}
		\tableofcontents
	\end{frame}
	
	% ============================================================
	\section{1. Enterprise Architecture}
	% ============================================================
	
	\begin{frame}{Q201: Architecture Decision Records}
		\begin{block}{Question}
			Why are Architecture Decision Records (ADRs) valuable for governance?
		\end{block}
		\begin{enumerate}[(A)]
			\item They document every meeting minute
			\item They capture context, rationale, and consequences of key architecture decisions
			\item They replace architecture reviews
			\item They eliminate risk
		\end{enumerate}
	\end{frame}
	
	\begin{frame}{Q201: Explanation}
		\begin{block}{Correct Answer: B}
			\textbf{ADRs preserve rationale:}
			\begin{itemize}
				\item Capture context and options considered
				\item Explain why a decision was made
				\item Support audit and future teams
			\end{itemize}
		\end{block}
		\begin{alertblock}{Critical Trap}
			ADRs record decisions and reasoning, not meetings.
		\end{alertblock}
	\end{frame}
	
	\begin{frame}{Q202: Capability-Based Planning}
		\begin{block}{Question}
			What does capability-based planning in EA emphasize?
		\end{block}
		\begin{enumerate}[(A)]
			\item Vendor preferences
			\item Delivering business capabilities that enable strategy
			\item Buying the newest tools
			\item Cutting costs only
		\end{enumerate}
	\end{frame}
	
	\begin{frame}{Q202: Explanation}
		\begin{block}{Correct Answer: B}
			\textbf{Capabilities enable strategy:}
			\begin{itemize}
				\item Focus on business capabilities
				\item Technology supports capabilities
				\item Reduces technology-driven decisions
			\end{itemize}
		\end{block}
		\begin{alertblock}{Critical Trap}
			Capability-based planning starts from business capabilities.
		\end{alertblock}
	\end{frame}
	
	\begin{frame}{Q203: Architecture Debt Trade-Off}
		\begin{block}{Question}
			When might it be ACCEPTABLE to take on architecture debt?
		\end{block}
		\begin{enumerate}[(A)]
			\item Always, to save time
			\item When the debt is time-bound and has a documented repayment plan
			\item Never
			\item When vendors ask
		\end{enumerate}
	\end{frame}
	
	\begin{frame}{Q203: Explanation}
		\begin{block}{Correct Answer: B}
			\textbf{Managed debt is acceptable:}
			\begin{itemize}
				\item Short-term business need
				\item Documented and tracked
				\item Repayment plan agreed
			\end{itemize}
		\end{block}
		\begin{alertblock}{Critical Trap}
			Architecture debt is acceptable only if managed and repaid.
		\end{alertblock}
	\end{frame}
	
	\begin{frame}{Q204: TOGAF Architecture Governance}
		\begin{block}{Question}
			What does TOGAF architecture governance typically include?
		\end{block}
		\begin{enumerate}[(A)]
			\item Only documentation
			\item Architecture board, compliance reviews, and dispensations
			\item Only marketing
			\item Only legal
		\end{enumerate}
	\end{frame}
	
	\begin{frame}{Q204: Explanation}
		\begin{block}{Correct Answer: B}
			\textbf{Governance includes board and dispensations:}
			\begin{itemize}
				\item Architecture board
				\item Compliance reviews
				\item Dispensations for exceptions
			\end{itemize}
		\end{block}
		\begin{alertblock}{Critical Trap}
			Governance includes exception mechanisms, not just rules.
		\end{alertblock}
	\end{frame}
	
	\begin{frame}{Q205: Business Architecture Artifacts}
		\begin{block}{Question}
			Which is a typical business architecture artifact?
		\end{block}
		\begin{enumerate}[(A)]
			\item Firewall rules
			\item Business process model and capability map
			\item Server configuration
			\item Network diagrams
		\end{enumerate}
	\end{frame}
	
	\begin{frame}{Q205: Explanation}
		\begin{block}{Correct Answer: B}
			\textbf{Business architecture uses process and capability models:}
			\begin{itemize}
				\item Business processes
				\item Capability maps
				\item Organizational models
			\end{itemize}
		\end{block}
		\begin{alertblock}{Critical Trap}
			Business architecture = processes and capabilities.
		\end{alertblock}
	\end{frame}
	
	\begin{frame}{Q206: Application Portfolio Management}
		\begin{block}{Question}
			What does application portfolio management aim to do?
		\end{block}
		\begin{enumerate}[(A)]
			\item Buy more apps
			\item Rationalize apps to remove redundancy and optimize value
			\item Ignore business needs
			\item Eliminate IT
		\end{enumerate}
	\end{frame}
	
	\begin{frame}{Q206: Explanation}
		\begin{block}{Correct Answer: B}
			\textbf{Rationalize and optimize:}
			\begin{itemize}
				\item Identify redundant apps
				\item Retire low-value apps
				\item Align portfolio to business
			\end{itemize}
		\end{block}
		\begin{alertblock}{Critical Trap}
			APM reduces duplication and improves value.
		\end{alertblock}
	\end{frame}
	
	\begin{frame}{Q207: Reference Model Use}
		\begin{block}{Question}
			What is a risk of relying on reference models without adaptation?
		\end{block}
		\begin{enumerate}[(A)]
			\item Guaranteed success
			\item Mismatch with actual business context and priorities
			\item Reduced cost
			\item Perfect fit
		\end{enumerate}
	\end{frame}
	
	\begin{frame}{Q207: Explanation}
		\begin{block}{Correct Answer: B}
			\textbf{Reference models need tailoring:}
			\begin{itemize}
				\item Context differs
				\item Priorities differ
				\item Adaptation is required
			\end{itemize}
		\end{block}
		\begin{alertblock}{Critical Trap}
			Reference models are starting points, not mandates.
		\end{alertblock}
	\end{frame}
	
	\begin{frame}{Q208: Architecture Requirements}
		\begin{block}{Question}
			What should architecture requirements be derived from?
		\end{block}
		\begin{enumerate}[(A)]
			\item Vendor marketing
			\item Business requirements and strategy
			\item Personal preferences
			\item Social media
		\end{enumerate}
	\end{frame}
	
	\begin{frame}{Q208: Explanation}
		\begin{block}{Correct Answer: B}
			\textbf{Architecture serves business:}
			\begin{itemize}
				\item Business requirements
				\item Strategy alignment
				\item Stakeholder needs
			\end{itemize}
		\end{block}
		\begin{alertblock}{Critical Trap}
			Architecture requirements follow business, not vendors.
		\end{alertblock}
	\end{frame}
	
	\begin{frame}{Q209: Architecture Review Value}
		\begin{block}{Question}
			What is the primary value of an architecture review?
		\end{block}
		\begin{enumerate}[(A)]
			\item Slowing down projects
			\item Detecting misalignment and risk before investment
			\item Increasing costs
			\item Avoiding governance
		\end{enumerate}
	\end{frame}
	
	\begin{frame}{Q209: Explanation}
		\begin{block}{Correct Answer: B}
			\textbf{Review catches issues early:}
			\begin{itemize}
				\item Detects misalignment
				\item Reduces rework
				\item Saves cost
			\end{itemize}
		\end{block}
		\begin{alertblock}{Critical Trap}
			Architecture reviews reduce risk, not add bureaucracy.
		\end{alertblock}
	\end{frame}
	
	\begin{frame}{Q210: EA Roadmap}
		\begin{block}{Question}
			What does an EA roadmap provide?
		\end{block}
		\begin{enumerate}[(A)]
			\item Vendor list
			\item Sequenced transition plan from current to target state
			\item Only budgets
			\item Only diagrams
		\end{enumerate}
	\end{frame}
	
	\begin{frame}{Q210: Explanation}
		\begin{block}{Correct Answer: B}
			\textbf{Roadmap sequences the transition:}
			\begin{itemize}
				\item Current to target state
				\item Prioritized projects
				\item Dependencies and milestones
			\end{itemize}
		\end{block}
		\begin{alertblock}{Critical Trap}
			Roadmaps sequence change, not just document it.
		\end{alertblock}
	\end{frame}
	
	% ============================================================
	\section{2. IT Operations and Change Management}
	% ============================================================
	
	\begin{frame}{Q211: Change Backout Plan}
		\begin{block}{Question}
			What should a backout plan include?
		\end{block}
		\begin{enumerate}[(A)]
			\item Only approval signatures
			\item Steps, verification, and communication to revert the change safely
			\item Only costs
			\item Only vendor names
		\end{enumerate}
	\end{frame}
	
	\begin{frame}{Q211: Explanation}
		\begin{block}{Correct Answer: B}
			\textbf{Backout plan must be executable:}
			\begin{itemize}
				\item Clear steps
				\item Verification methods
				\item Communication plan
			\end{itemize}
		\end{block}
		\begin{alertblock}{Critical Trap}
			Backout plan must be testable and clear.
		\end{alertblock}
	\end{frame}
	
	\begin{frame}{Q212: Problem Management Value}
		\begin{block}{Question}
			Why is problem management valuable for risk reduction?
		\end{block}
		\begin{enumerate}[(A)]
			\item Faster tickets
			\item It removes recurring incidents by addressing root causes
			\item It increases workload
			\item It avoids audits
		\end{enumerate}
	\end{frame}
	
	\begin{frame}{Q212: Explanation}
		\begin{block}{Correct Answer: B}
			\textbf{Problem management eliminates root causes:}
			\begin{itemize}
				\item Reduces recurring incidents
				\item Frees resources
				\item Improves stability
			\end{itemize}
		\end{block}
		\begin{alertblock}{Critical Trap}
			Problem management reduces risk by eliminating causes.
		\end{alertblock}
	\end{frame}
	
	\begin{frame}{Q213: Release Calendar}
		\begin{block}{Question}
			Why is a release calendar valuable?
		\end{block}
		\begin{enumerate}[(A)]
			\item Marketing
			\item Coordinates releases to avoid conflicts and business disruption
			\item Hides changes from business
			\item Ignores testing
		\end{enumerate}
	\end{frame}
	
	\begin{frame}{Q213: Explanation}
		\begin{block}{Correct Answer: B}
			\textbf{Calendar coordinates deployment:}
			\begin{itemize}
				\item Avoids release collisions
				\item Aligns with business periods
				\item Reduces risk
			\end{itemize}
		\end{block}
		\begin{alertblock}{Critical Trap}
			Release calendar prevents chaos, not just schedules.
		\end{alertblock}
	\end{frame}
	
	\begin{frame}{Q214: Test Environment Parity}
		\begin{block}{Question}
			Why should test environments mirror production closely?
		\end{block}
		\begin{enumerate}[(A)]
			\item Cheaper
			\item To catch issues that would occur in production before release
			\item Faster
			\item Smaller
		\end{enumerate}
	\end{frame}
	
	\begin{frame}{Q214: Explanation}
		\begin{block}{Correct Answer: B}
			\textbf{Parity catches real issues:}
			\begin{itemize}
				\item Configuration and dependency differences cause failures
				\item Production parity reduces surprises
				\item Testing quality improves
			\end{itemize}
		\end{block}
		\begin{alertblock}{Critical Trap}
			Unrealistic test environments miss real issues.
		\end{alertblock}
	\end{frame}
	
	\begin{frame}{Q215: Change Emergency Retrospective}
		\begin{block}{Question}
			What should happen after an emergency change?
		\end{block}
		\begin{enumerate}[(A)]
			\item Nothing
			\item Retrospective review and document to verify controls and learn
			\item Cancel it
			\item Ignore it
		\end{enumerate}
	\end{frame}
	
	\begin{frame}{Q215: Explanation}
		\begin{block}{Correct Answer: B}
			\textbf{Emergency changes require retrospective:}
			\begin{itemize}
				\item Verify it was justified
				\item Confirm testing occurred
				\item Capture lessons
			\end{itemize}
		\end{block}
		\begin{alertblock}{Critical Trap}
			Emergency without retrospective is unmanaged.
		\end{alertblock}
	\end{frame}
	
	\begin{frame}{Q216: Runbook Value}
		\begin{block}{Question}
			What is the value of an operations runbook?
		\end{block}
		\begin{enumerate}[(A)]
			\item Marketing
			\item Consistent execution of operational tasks and recovery procedures
			\item Higher cost
			\item Larger teams
		\end{enumerate}
	\end{frame}
	
	\begin{frame}{Q216: Explanation}
		\begin{block}{Correct Answer: B}
			\textbf{Runbooks standardize execution:}
			\begin{itemize}
				\item Reduce errors
				\item Consistent response
				\item Faster onboarding
			\end{itemize}
		\end{block}
		\begin{alertblock}{Critical Trap}
			Runbooks reduce reliance on memory.
		\end{alertblock}
	\end{frame}
	
	\begin{frame}{Q217: Operations KPI Choice}
		\begin{block}{Question}
			Which KPI is most useful for operations leadership?
		\end{block}
		\begin{enumerate}[(A)]
			\item Number of tickets opened
			\item Availability and mean time to resolve critical incidents
			\item Number of meetings
			\item Number of employees
		\end{enumerate}
	\end{frame}
	
	\begin{frame}{Q217: Explanation}
		\begin{block}{Correct Answer: B}
			\textbf{Availability and MTTR show stability:}
			\begin{itemize}
				\item Availability impacts business
				\item MTTR shows response capability
				\item Ticket volume is workload, not outcome
			\end{itemize}
		\end{block}
		\begin{alertblock}{Critical Trap}
			Leadership cares about stability, not workload.
		\end{alertblock}
	\end{frame}
	
	\begin{frame}{Q218: Monitoring Coverage}
		\begin{block}{Question}
			Why is monitoring coverage important?
		\end{block}
		\begin{enumerate}[(A)]
			\item Cheaper
			\item Detects issues proactively before users report them
			\item Faster tickets
			\item Less training
		\end{enumerate}
	\end{frame}
	
	\begin{frame}{Q218: Explanation}
		\begin{block}{Correct Answer: B}
			\textbf{Coverage prevents surprises:}
			\begin{itemize}
				\item Proactive detection
				\item Reduced downtime
				\item Better user experience
			\end{itemize}
		\end{block}
		\begin{alertblock}{Critical Trap}
			Unmonitored systems fail silently.
		\end{alertblock}
	\end{frame}
	
	\begin{frame}{Q219: Alerting Tuning}
		\begin{block}{Question}
			What is the main benefit of tuning alerts?
		\end{block}
		\begin{enumerate}[(A)]
			\item Fewer alerts always means safer
			\item Reduces noise so critical alerts stand out and response improves
			\item Cheaper
			\item Slower response
		\end{enumerate}
	\end{frame}
	
	\begin{frame}{Q219: Explanation}
		\begin{block}{Correct Answer: B}
			\textbf{Tuning improves signal:}
			\begin{itemize}
				\item Noise hides critical alerts
				\item Tuning improves response
				\item Alert fatigue is a risk
			\end{itemize}
		\end{block}
		\begin{alertblock}{Critical Trap}
			Too many alerts means too few responses.
		\end{alertblock}
	\end{frame}
	
	\begin{frame}{Q220: Post-Incident Actions}
		\begin{block}{Question}
			Which is the MOST important output of an incident review?
		\end{block}
		\begin{enumerate}[(A)]
			\item Blame
			\item Action items with owners and timelines to reduce recurrence
			\item Ticket closure
			\item Budget increase
		\end{enumerate}
	\end{frame}
	
	\begin{frame}{Q220: Explanation}
		\begin{block}{Correct Answer: B}
			\textbf{Actions drive improvement:}
			\begin{itemize}
				\item Actions drive improvement
				\item Owners ensure follow-through
				\item Timelines create accountability
			\end{itemize}
		\end{block}
		\begin{alertblock}{Critical Trap}
			Incident reviews without actions are theater.
		\end{alertblock}
	\end{frame}
	
	% ============================================================
	\section{3. Enterprise Resiliency}
	% ============================================================
	
	\begin{frame}{Q221: Recovery Site Selection}
		\begin{block}{Question}
			Which factor MOST influences recovery site selection?
		\end{block}
		\begin{enumerate}[(A)]
			\item Cost only
			\item Business RTO, RPO, and cost trade-offs
			\item Vendor preference
			\item Marketing
		\end{enumerate}
	\end{frame}
	
	\begin{frame}{Q221: Explanation}
		\begin{block}{Correct Answer: B}
			\textbf{Balance recovery with cost:}
			\begin{itemize}
				\item RTO and RPO define requirements
				\item Cost constrains choices
				\item Site fits business needs
			\end{itemize}
		\end{block}
		\begin{alertblock}{Critical Trap}
			Site selection balances recovery objectives with cost.
		\end{alertblock}
	\end{frame}
	
	\begin{frame}{Q222: Cloud DR Benefits}
		\begin{block}{Question}
			What is the MAIN benefit of cloud-based DR?
		\end{block}
		\begin{enumerate}[(A)]
			\item Free
			\item Scalable recovery without owning a second data center
			\item No testing needed
			\item No security controls
		\end{enumerate}
	\end{frame}
	
	\begin{frame}{Q222: Explanation}
		\begin{block}{Correct Answer: B}
			\textbf{Scale without owning a second site:}
			\begin{itemize}
				\item Rapid scaling
				\item Reduced cost of standby
				\item Still requires testing and security
			\end{itemize}
		\end{block}
		\begin{alertblock}{Critical Trap}
			Cloud DR still requires testing and security.
		\end{alertblock}
	\end{frame}
	
	\begin{frame}{Q223: Failover Automation Risk}
		\begin{block}{Question}
			What risk is introduced by fully automated failover?
		\end{block}
		\begin{enumerate}[(A)]
			\item Slower recovery
			\item Automated failover can be triggered by false positives and cause outage
			\item Higher cost
			\item Fewer tests
		\end{enumerate}
	\end{frame}
	
	\begin{frame}{Q223: Explanation}
		\begin{block}{Correct Answer: B}
			\textbf{False positives cause outages:}
			\begin{itemize}
				\item False triggers cause disruption
				\item Requires careful tuning
				\item Manual override should exist
			\end{itemize}
		\end{block}
		\begin{alertblock}{Critical Trap}
			Automation reduces time but can misfire.
		\end{alertblock}
	\end{frame}
	
	\begin{frame}{Q224: Backup Encryption}
		\begin{block}{Question}
			Why should backups be encrypted?
		\end{block}
		\begin{enumerate}[(A)]
			\item Slower restore
			\item Protect sensitive data if backups are lost or stolen
			\item Cheaper
			\item Less space
		\end{enumerate}
	\end{frame}
	
	\begin{frame}{Q224: Explanation}
		\begin{block}{Correct Answer: B}
			\textbf{Protect sensitive backup data:}
			\begin{itemize}
				\item Backups contain sensitive data
				\item Encryption protects if media lost
				\item Keys must be protected separately
			\end{itemize}
		\end{block}
		\begin{alertblock}{Critical Trap}
			Unencrypted backups are a data exposure risk.
		\end{alertblock}
	\end{frame}
	
	\begin{frame}{Q225: Backup Isolation}
		\begin{block}{Question}
			What is the purpose of isolating backups from production networks?
		\end{block}
		\begin{enumerate}[(A)]
			\item Speed
			\item Prevent attackers from reaching and encrypting backups
			\item Cheaper
			\item Easier restore
		\end{enumerate}
	\end{frame}
	
	\begin{frame}{Q225: Explanation}
		\begin{block}{Correct Answer: B}
			\textbf{Isolation protects backups:}
			\begin{itemize}
				\item Ransomware spreads
				\item Isolation protects backups
				\item Offline or air gap best
			\end{itemize}
		\end{block}
		\begin{alertblock}{Critical Trap}
			Connected backups are vulnerable to ransomware.
		\end{alertblock}
	\end{frame}
	
	\begin{frame}{Q226: BCP Coverage}
		\begin{block}{Question}
			What should a BCP cover beyond IT recovery?
		\end{block}
		\begin{enumerate}[(A)]
			\item Only servers
			\item People, facilities, suppliers, communications, and processes
			\item Only backups
			\item Only firewalls
		\end{enumerate}
	\end{frame}
	
	\begin{frame}{Q226: Explanation}
		\begin{block}{Correct Answer: B}
			\textbf{Continuity is business-wide:}
			\begin{itemize}
				\item Business continuity is broad
				\item People and facilities matter
				\item Suppliers matter
				\item Communications matter
			\end{itemize}
		\end{block}
		\begin{alertblock}{Critical Trap}
			BCP covers more than IT; it covers the business.
		\end{alertblock}
	\end{frame}
	
	\begin{frame}{Q227: Crisis Leadership}
		\begin{block}{Question}
			What is MOST important in crisis leadership?
		\end{block}
		\begin{enumerate}[(A)]
			\item Hiding the problem
			\item Clear decisions, communication, and direction with delegated authority
			\item Slower meetings
			\item More approvals
		\end{enumerate}
	\end{frame}
	
	\begin{frame}{Q227: Explanation}
		\begin{block}{Correct Answer: B}
			\textbf{Clear decisions and communication:}
			\begin{itemize}
				\item Crises require speed
				\item Clear authority
				\item Consistent communication
			\end{itemize}
		\end{block}
		\begin{alertblock}{Critical Trap}
			Ambiguous leadership lengthens crises.
		\end{alertblock}
	\end{frame}
	
	\begin{frame}{Q228: Vendor Dependency in DR}
		\begin{block}{Question}
			What risk exists if a DR provider is also a single point of failure?
		\end{block}
		\begin{enumerate}[(A)]
			\item No risk
			\item Concentration risk if the provider fails during the same event
			\item Cheaper recovery
			\item Faster recovery
		\end{enumerate}
	\end{frame}
	
	\begin{frame}{Q228: Explanation}
		\begin{block}{Correct Answer: B}
			\textbf{Concentrated DR is risky:}
			\begin{itemize}
				\item Shared geography or vendor
				\item Concentrated risk
				\item Use diverse providers
			\end{itemize}
		\end{block}
		\begin{alertblock}{Critical Trap}
			DR providers can be part of a concentration risk.
		\end{alertblock}
	\end{frame}
	
	\begin{frame}{Q229: Recovery Metrics}
		\begin{block}{Question}
			Which metric measures data loss tolerance?
		\end{block}
		\begin{enumerate}[(A)]
			\item RTO
			\item RPO
			\item MTD
			\item MTTR
		\end{enumerate}
	\end{frame}
	
	\begin{frame}{Q229: Explanation}
		\begin{block}{Correct Answer: B}
			\textbf{RPO measures data loss:}
			\begin{itemize}
				\item RPO is acceptable data loss
				\item RTO is recovery time target
				\item MTD is max tolerable outage
			\end{itemize}
		\end{block}
		\begin{alertblock}{Critical Trap}
			RPO is about data loss; RTO is about time.
		\end{alertblock}
	\end{frame}
	
	\begin{frame}{Q230: Third-Party Recovery Testing}
		\begin{block}{Question}
			Why is testing recovery for third-party dependencies important?
		\end{block}
		\begin{enumerate}[(A)]
			\item Trust contracts
			\item Verify recovery actually works with external providers, not just assumptions
			\item Cheaper
			\item Faster
		\end{enumerate}
	\end{frame}
	
	\begin{frame}{Q230: Explanation}
		\begin{block}{Correct Answer: B}
			\textbf{Tests verify promises:}
			\begin{itemize}
				\item Contracts promise
				\item Tests verify
				\item Real world differs
			\end{itemize}
		\end{block}
		\begin{alertblock}{Critical Trap}
			Contracts do not replace recovery testing.
		\end{alertblock}
	\end{frame}
	
	% ============================================================
	\section{4. Data Lifecycle Management}
	% ============================================================
	
	\begin{frame}{Q231: Data Lineage}
		\begin{block}{Question}
			What is the purpose of data lineage in data management?
		\end{block}
		\begin{enumerate}[(A)]
			\item Marketing
			\item Trace where data came from, how it moved, and where it is used
			\item Cheaper
			\item Backups
		\end{enumerate}
	\end{frame}
	
	\begin{frame}{Q231: Explanation}
		\begin{block}{Correct Answer: B}
			\textbf{Lineage provides traceability:}
			\begin{itemize}
				\item Traceability
				\item Supports governance
				\item Enables impact analysis
			\end{itemize}
		\end{block}
		\begin{alertblock}{Critical Trap}
			Lineage provides data traceability, not just storage.
		\end{alertblock}
	\end{frame}
	
	\begin{frame}{Q232: Data Masking}
		\begin{block}{Question}
			When should data masking be used?
		\end{block}
		\begin{enumerate}[(A)]
			\item Never
			\item When using production-like data in non-production environments
			\item Only in production
			\item Only for archives
		\end{enumerate}
	\end{frame}
	
	\begin{frame}{Q232: Explanation}
		\begin{block}{Correct Answer: B}
			\textbf{Masking for non-production:}
			\begin{itemize}
				\item Protects sensitive fields
				\item Enables realistic testing
				\item Avoids exposure
			\end{itemize}
		\end{block}
		\begin{alertblock}{Critical Trap}
			Masking is required when production-like data is used in testing.
		\end{alertblock}
	\end{frame}
	
	\begin{frame}{Q233: Pseudonymization vs Anonymization}
		\begin{block}{Question}
			What is the key difference?
		\end{block}
		\begin{enumerate}[(A)]
			\item They are the same
			\item Pseudonymization is reversible with a key; anonymization is not reversible
			\item Anonymization is reversible
			\item Pseudonymization is not reversible
		\end{enumerate}
	\end{frame}
	
	\begin{frame}{Q233: Explanation}
		\begin{block}{Correct Answer: B}
			\textbf{Reversibility differs:}
			\begin{itemize}
				\item Pseudonymization is reversible with key
				\item Anonymization is permanent
				\item GDPR treats them differently
			\end{itemize}
		\end{block}
		\begin{alertblock}{Critical Trap}
			Pseudonymized data is still personal data.
		\end{alertblock}
	\end{frame}
	
	\begin{frame}{Q234: Data Retention Audit}
		\begin{block}{Question}
			Why should retention schedules be audited?
		\end{block}
		\begin{enumerate}[(A)]
			\item To show effort
			\item To verify data is not retained beyond legal limits
			\item To reduce storage only
			\item To avoid training
		\end{enumerate}
	\end{frame}
	
	\begin{frame}{Q234: Explanation}
		\begin{block}{Correct Answer: B}
			\textbf{Audit ensures compliance:}
			\begin{itemize}
				\item Over-retention risk
				\item Legal exposure
				\item Compliance requires proof
			\end{itemize}
		\end{block}
		\begin{alertblock}{Critical Trap}
			Retention schedules must be enforced and audited.
		\end{alertblock}
	\end{frame}
	
	\begin{frame}{Q235: Data Steward Role}
		\begin{block}{Question}
			What does a data steward typically do?
		\end{block}
		\begin{enumerate}[(A)]
			\item Runs the helpdesk
			\item Maintains quality, classification, and metadata for data under their care
			\item Owns the firewall
			\item Manages HR
		\end{enumerate}
	\end{frame}
	
	\begin{frame}{Q235: Explanation}
		\begin{block}{Correct Answer: B}
			\textbf{Stewards manage quality and metadata:}
			\begin{itemize}
				\item Quality assurance
				\item Metadata
				\item Compliance with data policy
			\end{itemize}
		\end{block}
		\begin{alertblock}{Critical Trap}
			Stewards manage data quality and metadata, not systems.
		\end{alertblock}
	\end{frame}
	
	\begin{frame}{Q236: Data Classification Oversharing}
		\begin{block}{Question}
			What is the risk of over-classifying data as highly confidential?
		\end{block}
		\begin{enumerate}[(A)]
			\item Improved productivity
			\item Excessive controls slow work, cost more, and reduce trust in classification
			\item Higher productivity
			\item Lower cost
		\end{enumerate}
	\end{frame}
	
	\begin{frame}{Q236: Explanation}
		\begin{block}{Correct Answer: B}
			\textbf{Over-classification is costly:}
			\begin{itemize}
				\item Controls proportionate
				\item Over-classification erodes compliance
				\item Classification must be accurate
			\end{itemize}
		\end{block}
		\begin{alertblock}{Critical Trap}
			Over-classification is as damaging as under-classification.
		\end{alertblock}
	\end{frame}
	
	\begin{frame}{Q237: Cross-Border Retention}
		\begin{block}{Question}
			Why do data retention rules differ across jurisdictions?
		\end{block}
		\begin{enumerate}[(A)]
			\item No reason
			\item Different laws, cultures, and regulatory priorities
			\item Random
			\item Marketing
		\end{enumerate}
	\end{frame}
	
	\begin{frame}{Q237: Explanation}
		\begin{block}{Correct Answer: B}
			\textbf{Legal variation:}
			\begin{itemize}
				\item Legal variation
				\item Sector-specific rules
				\item Multi-jurisdiction complexity
			\end{itemize}
		\end{block}
		\begin{alertblock}{Critical Trap}
			Multi-jurisdiction retention requires legal mapping.
		\end{alertblock}
	\end{frame}
	
	\begin{frame}{Q238: Data Destruction Verification}
		\begin{block}{Question}
			How should an organization verify destruction on sanitized media?
		\end{block}
		\begin{enumerate}[(A)]
			\item Trust the vendor
			\item Certificates, third-party verification, or sampling
			\item Assume
			\item Ignore
		\end{enumerate}
	\end{frame}
	
	\begin{frame}{Q238: Explanation}
		\begin{block}{Correct Answer: B}
			\textbf{Verify destruction:}
			\begin{itemize}
				\item Certificates
				\item Sampling
				\item Audit trails
			\end{itemize}
		\end{block}
		\begin{alertblock}{Critical Trap}
			Destruction without verification is not defensible.
		\end{alertblock}
	\end{frame}
	
	\begin{frame}{Q239: DLP Coverage Limits}
		\begin{block}{Question}
			What is a limitation of DLP tools?
		\end{block}
		\begin{enumerate}[(A)]
			\item Perfect detection
			\item Cannot catch all sensitive data, especially unstructured content or encrypted channels
			\item Cheaper than alternatives
			\item Replaces training
		\end{enumerate}
	\end{frame}
	
	\begin{frame}{Q239: Explanation}
		\begin{block}{Correct Answer: B}
			\textbf{DLP has limits:}
			\begin{itemize}
				\item Content inspection limits
				\item Encrypted traffic bypasses
				\item Complement with other controls
			\end{itemize}
		\end{block}
		\begin{alertblock}{Critical Trap}
			DLP is not a silver bullet; combine controls.
		\end{alertblock}
	\end{frame}
	
	\begin{frame}{Q240: Data Privacy by Default}
		\begin{block}{Question}
			What does privacy by default require in data lifecycle?
		\end{block}
		\begin{enumerate}[(A)]
			\item Collect all data
			\item Default settings minimize data collection and sharing
			\item Force users to opt out
			\item Ignore regulations
		\end{enumerate}
	\end{frame}
	
	\begin{frame}{Q240: Explanation}
		\begin{block}{Correct Answer: B}
			\textbf{Minimize by default:}
			\begin{itemize}
				\item Strictest settings default
				\item Users opt in to sharing
				\item GDPR principle
			\end{itemize}
		\end{block}
		\begin{alertblock}{Critical Trap}
			Privacy by default = minimized data sharing by default.
		\end{alertblock}
	\end{frame}
	
	% ============================================================
	\section{5. SDLC and Emerging Technologies}
	% ============================================================
	
	\begin{frame}{Q241: DevSecOps Principle}
		\begin{block}{Question}
			What is the core principle of DevSecOps?
		\end{block}
		\begin{enumerate}[(A)]
			\item Security after deployment
			\item Security integrated throughout development and operations
			\item Security only in audit
			\item No security
		\end{enumerate}
	\end{frame}
	
	\begin{frame}{Q241: Explanation}
		\begin{block}{Correct Answer: B}
			\textbf{Security in every phase:}
			\begin{itemize}
				\item Security in every phase
				\item Automation of checks
				\item Shared responsibility
			\end{itemize}
		\end{block}
		\begin{alertblock}{Critical Trap}
			DevSecOps integrates security, it does not add it late.
		\end{alertblock}
	\end{frame}
	
	\begin{frame}{Q242: CI/CD Security Gates}
		\begin{block}{Question}
			What role do security gates play in CI/CD?
		\end{block}
		\begin{enumerate}[(A)]
			\item Slow down deployments only
			\item Prevent vulnerable or non-compliant code from reaching production
			\item Cheaper
			\item Skip tests
		\end{enumerate}
	\end{frame}
	
	\begin{frame}{Q242: Explanation}
		\begin{block}{Correct Answer: B}
			\textbf{Gates block risk:}
			\begin{itemize}
				\item SAST, DAST, dependency scans
				\item Block on critical issues
				\item Automate policy
			\end{itemize}
		\end{block}
		\begin{alertblock}{Critical Trap}
			Security gates block risk, not velocity.
		\end{alertblock}
	\end{frame}
	
	\begin{frame}{Q243: Infrastructure as Code Risk}
		\begin{block}{Question}
			What is the biggest risk in Infrastructure as Code?
		\end{block}
		\begin{enumerate}[(A)]
			\item Automation
			\item Misconfiguration propagated at scale
			\item Cheaper
			\item Faster
		\end{enumerate}
	\end{frame}
	
	\begin{frame}{Q243: Explanation}
		\begin{block}{Correct Answer: B}
			\textbf{Errors scale with IaC:}
			\begin{itemize}
				\item Errors scale
				\item Secrets in code
				\item Review and scanning essential
			\end{itemize}
		\end{block}
		\begin{alertblock}{Critical Trap}
			IaC errors replicate quickly; review and scan.
		\end{alertblock}
	\end{frame}
	
	\begin{frame}{Q244: Secret Management}
		\begin{block}{Question}
			Why should secrets be stored in a secrets manager?
		\end{block}
		\begin{enumerate}[(A)]
			\item Cheaper
			\item Centralize protection, rotation, and access control of secrets
			\item Slower
			\item Easier to leak
		\end{enumerate}
	\end{frame}
	
	\begin{frame}{Q244: Explanation}
		\begin{block}{Correct Answer: B}
			\textbf{Secrets managers centralize control:}
			\begin{itemize}
				\item Central control
				\item Rotation
				\item Audit access
			\end{itemize}
		\end{block}
		\begin{alertblock}{Critical Trap}
			Hardcoded secrets are a breach waiting to happen.
		\end{alertblock}
	\end{frame}
	
	\begin{frame}{Q245: Serverless Security}
		\begin{block}{Question}
			What is a key security consideration in serverless computing?
		\end{block}
		\begin{enumerate}[(A)]
			\item No security needed
			\item Identity, permissions, and dependency risks in functions
			\item Only backups
			\item Only firewalls
		\end{enumerate}
	\end{frame}
	
	\begin{frame}{Q245: Explanation}
		\begin{block}{Correct Answer: B}
			\textbf{Serverless still needs controls:}
			\begin{itemize}
				\item Fine-grained IAM
				\item Dependency scanning
				\item Logging
			\end{itemize}
		\end{block}
		\begin{alertblock}{Critical Trap}
			Serverless still requires IAM and dependency controls.
		\end{alertblock}
	\end{frame}
	
	\begin{frame}{Q246: API Security}
		\begin{block}{Question}
			Which is a top risk for APIs?
		\end{block}
		\begin{enumerate}[(A)]
			\item Marketing
			\item Broken authentication and authorization
			\item Cheaper
			\item Backup cost
		\end{enumerate}
	\end{frame}
	
	\begin{frame}{Q246: Explanation}
		\begin{block}{Correct Answer: B}
			\textbf{Auth flaws are top API risks:}
			\begin{itemize}
				\item AuthN and AuthZ issues
				\item Rate limiting
				\item Input validation
			\end{itemize}
		\end{block}
		\begin{alertblock}{Critical Trap}
			APIs are often attacked via auth flaws.
		\end{alertblock}
	\end{frame}
	
	\begin{frame}{Q247: Mobile App Security}
		\begin{block}{Question}
			Which is a key risk for mobile apps?
		\end{block}
		\begin{enumerate}[(A)]
			\item Marketing
			\item Insecure storage of data on devices
			\item Cheaper devices
			\item Screen size
		\end{enumerate}
	\end{frame}
	
	\begin{frame}{Q247: Explanation}
		\begin{block}{Correct Answer: B}
			\textbf{Local storage risks:}
			\begin{itemize}
				\item Local storage risks
				\item Keychain or keystore use
				\item Encryption at rest
			\end{itemize}
		\end{block}
		\begin{alertblock}{Critical Trap}
			Mobile devices are lost and stolen; protect local data.
		\end{alertblock}
	\end{frame}
	
	\begin{frame}{Q248: Zero-Trust Microsegmentation}
		\begin{block}{Question}
			What is microsegmentation?
		\end{block}
		\begin{enumerate}[(A)]
			\item Divide networks into very small, isolated zones
			\item Backup strategy
			\item Marketing term
			\item Hiring practice
		\end{enumerate}
	\end{frame}
	
	\begin{frame}{Q248: Explanation}
		\begin{block}{Correct Answer: A}
			\textbf{Microsegmentation isolates zones:}
			\begin{itemize}
				\item Fine-grained segmentation
				\item Limits lateral movement
				\item Common in Zero Trust
			\end{itemize}
		\end{block}
		\begin{alertblock}{Critical Trap}
			Microsegmentation reduces blast radius.
		\end{alertblock}
	\end{frame}
	
	\begin{frame}{Q249: Post-Quantum Migration}
		\begin{block}{Question}
			What is the recommended approach for post-quantum readiness?
		\end{block}
		\begin{enumerate}[(A)]
			\item Wait until quantum computers are practical
			\item Inventory cryptography, plan migration, adopt hybrid approaches
			\item Ignore
			\item Replace all systems immediately
		\end{enumerate}
	\end{frame}
	
	\begin{frame}{Q249: Explanation}
		\begin{block}{Correct Answer: B}
			\textbf{Plan post-quantum migration:}
			\begin{itemize}
				\item Inventory current crypto
				\item Plan migration
				\item Hybrid schemes
			\end{itemize}
		\end{block}
		\begin{alertblock}{Critical Trap}
			Quantum migration must be planned years in advance.
		\end{alertblock}
	\end{frame}
	
	\begin{frame}{Q250: AI Data Poisoning}
		\begin{block}{Question}
			What is data poisoning in AI systems?
		\end{block}
		\begin{enumerate}[(A)]
			\item Healthy training
			\item Manipulating training data to corrupt model behavior
			\item Faster training
			\item Cheaper
		\end{enumerate}
	\end{frame}
	
	\begin{frame}{Q250: Explanation}
		\begin{block}{Correct Answer: B}
			\textbf{Poisoned data corrupts models:}
			\begin{itemize}
				\item Attacks training data
				\item Biases or corrupts outcomes
				\item Data validation matters
			\end{itemize}
		\end{block}
		\begin{alertblock}{Critical Trap}
			Poisoned training data corrupts model outputs.
		\end{alertblock}
	\end{frame}
	
	% ============================================================
	\section{6. Information Security Principles}
	% ============================================================
	
	\begin{frame}{Q251: Non-Repudiation Principle}
		\begin{block}{Question}
			Which control BEST supports non-repudiation?
		\end{block}
		\begin{enumerate}[(A)]
			\item Firewall
			\item Digital signatures and immutable audit logs
			\item Antivirus
			\item Backups
		\end{enumerate}
	\end{frame}
	
	\begin{frame}{Q251: Explanation}
		\begin{block}{Correct Answer: B}
			\textbf{Signatures and logs support non-repudiation:}
			\begin{itemize}
				\item Signatures prove origin
				\item Immutable logs preserve evidence
				\item Both support non-repudiation
			\end{itemize}
		\end{block}
		\begin{alertblock}{Critical Trap}
			Non-repudiation requires proof, not just policy.
		\end{alertblock}
	\end{frame}
	
	\begin{frame}{Q252: Encryption Reversibility}
		\begin{block}{Question}
			What distinguishes encryption from hashing?
		\end{block}
		\begin{enumerate}[(A)]
			\item Both are reversible
			\item Encryption is reversible with a key; hashing is one-way
			\item Both are one-way
			\item Neither is reversible
		\end{enumerate}
	\end{frame}
	
	\begin{frame}{Q252: Explanation}
		\begin{block}{Correct Answer: B}
			\textbf{Different purposes:}
			\begin{itemize}
				\item Encryption protects confidentiality
				\item Hashing protects integrity
				\item Different purposes
			\end{itemize}
		\end{block}
		\begin{alertblock}{Critical Trap}
			Encryption is reversible; hashing is not.
		\end{alertblock}
	\end{frame}
	
	\begin{frame}{Q253: Certificate Expiry}
		\begin{block}{Question}
			What happens when a certificate expires?
		\end{block}
		\begin{enumerate}[(A)]
			\item It is auto-renewed
			\item It is no longer valid and trusting parties reject it
			\item It is renewed by the browser
			\item It stays valid
		\end{enumerate}
	\end{frame}
	
	\begin{frame}{Q253: Explanation}
		\begin{block}{Correct Answer: B}
			\textbf{Expired certs are rejected:}
			\begin{itemize}
				\item No longer trusted
				\item Causes outages
				\item Monitor expiry
			\end{itemize}
		\end{block}
		\begin{alertblock}{Critical Trap}
			Expired certs break services.
		\end{alertblock}
	\end{frame}
	
	\begin{frame}{Q254: Symmetric Key Distribution}
		\begin{block}{Question}
			What is the MAIN challenge with symmetric encryption?
		\end{block}
		\begin{enumerate}[(A)]
			\item Slower speed
			\item Secure key distribution between parties
			\item Larger key sizes
			\item Longer ciphertexts
		\end{enumerate}
	\end{frame}
	
	\begin{frame}{Q254: Explanation}
		\begin{block}{Correct Answer: B}
			\textbf{Key distribution is the challenge:}
			\begin{itemize}
				\item Shared key must be distributed
				\item Asymmetric helps establish
				\item Hybrid encryption is common
			\end{itemize}
		\end{block}
		\begin{alertblock}{Critical Trap}
			Symmetric challenges are about key distribution, not speed.
		\end{alertblock}
	\end{frame}
	
	\begin{frame}{Q255: Asymmetric Performance}
		\begin{block}{Question}
			Why is asymmetric encryption rarely used for bulk data?
		\end{block}
		\begin{enumerate}[(A)]
			\item Not secure
			\item Slower and computationally expensive
			\item Keys are short
			\item No integrity
		\end{enumerate}
	\end{frame}
	
	\begin{frame}{Q255: Explanation}
		\begin{block}{Correct Answer: B}
			\textbf{Asymmetric is slow:}
			\begin{itemize}
				\item Asymmetric is slow
				\item Used for key exchange or signatures
				\item Hybrid for bulk
			\end{itemize}
		\end{block}
		\begin{alertblock}{Critical Trap}
			Asymmetric protects keys; symmetric protects bulk data.
		\end{alertblock}
	\end{frame}
	
	\begin{frame}{Q256: Digital Certificate Chain}
		\begin{block}{Question}
			What does certificate chaining provide?
		\end{block}
		\begin{enumerate}[(A)]
			\item Redundancy
			\item Trust anchor linking leaf certs to a trusted root CA
			\item Faster connections
			\item Cheaper
		\end{enumerate}
	\end{frame}
	
	\begin{frame}{Q256: Explanation}
		\begin{block}{Correct Answer: B}
			\textbf{Chain to trusted root:}
			\begin{itemize}
				\item Chain to trusted root
				\item Enables validation
				\item Root of trust matters
			\end{itemize}
		\end{block}
		\begin{alertblock}{Critical Trap}
			The trust chain ends at the root CA.
		\end{alertblock}
	\end{frame}
	
	\begin{frame}{Q257: HSM Use Case}
		\begin{block}{Question}
			Which is a typical use case for an HSM?
		\end{block}
		\begin{enumerate}[(A)]
			\item Marketing analytics
			\item Root CA key protection and payment key management
			\item Backup storage
			\item Email filtering
		\end{enumerate}
	\end{frame}
	
	\begin{frame}{Q257: Explanation}
		\begin{block}{Correct Answer: B}
			\textbf{HSMs protect high-value keys:}
			\begin{itemize}
				\item High-value keys
				\item Regulated environments
				\item Physical tamper resistance
			\end{itemize}
		\end{block}
		\begin{alertblock}{Critical Trap}
			HSMs protect high-value keys, not general storage.
		\end{alertblock}
	\end{frame}
	
	\begin{frame}{Q258: Access Recertification Gap}
		\begin{block}{Question}
			What is a common risk if access recertifications are skipped?
		\end{block}
		\begin{enumerate}[(A)]
			\item Nothing
			\item Access creep and unauthorized access accumulating
			\item Faster access
			\item Fewer users
		\end{enumerate}
	\end{frame}
	
	\begin{frame}{Q258: Explanation}
		\begin{block}{Correct Answer: B}
			\textbf{Access creep:}
			\begin{itemize}
				\item Access accumulates
				\item Ex-employees retain access
				\item Privilege creep
			\end{itemize}
		\end{block}
		\begin{alertblock}{Critical Trap}
			Skipping recertification enables access creep.
		\end{alertblock}
	\end{frame}
	
	\begin{frame}{Q259: Behavior Analytics}
		\begin{block}{Question}
			What is user and entity behavior analytics primarily used for?
		\end{block}
		\begin{enumerate}[(A)]
			\item Marketing
			\item Detecting anomalous behavior that may indicate compromise
			\item Backup storage
			\item Encryption
		\end{enumerate}
	\end{frame}
	
	\begin{frame}{Q259: Explanation}
		\begin{block}{Correct Answer: B}
			\textbf{UEBA detects anomalies:}
			\begin{itemize}
				\item Baselines behavior
				\item Detects anomalies
				\item Complements SIEM
			\end{itemize}
		\end{block}
		\begin{alertblock}{Critical Trap}
			UEBA detects anomalies, not signatures only.
		\end{alertblock}
	\end{frame}
	
	\begin{frame}{Q260: Identity Provider Federation}
		\begin{block}{Question}
			What is the primary benefit of identity federation?
		\end{block}
		\begin{enumerate}[(A)]
			\item Marketing
			\item Users can access multiple systems with one trusted identity
			\item Cheaper
			\item Fewer controls
		\end{enumerate}
	\end{frame}
	
	\begin{frame}{Q260: Explanation}
		\begin{block}{Correct Answer: B}
			\textbf{SSO across domains:}
			\begin{itemize}
				\item Single identity
				\item Trust relationships
				\item SSO across domains
			\end{itemize}
		\end{block}
		\begin{alertblock}{Critical Trap}
			Federation enables SSO but requires trust management.
		\end{alertblock}
	\end{frame}
	
	\begin{frame}{Q261: Security Monitoring Maturity}
		\begin{block}{Question}
			Which factor MOST improves security monitoring maturity?
		\end{block}
		\begin{enumerate}[(A)]
			\item More alerts
			\item Tuned detection, playbooks, and skilled analysts
			\item More tools
			\item Larger budget
		\end{enumerate}
	\end{frame}
	
	\begin{frame}{Q261: Explanation}
		\begin{block}{Correct Answer: B}
			\textbf{People, tuning, and playbooks:}
			\begin{itemize}
				\item Tuning improves signal
				\item Playbooks speed response
				\item Skilled people matter
			\end{itemize}
		\end{block}
		\begin{alertblock}{Critical Trap}
			Tools alone do not improve maturity.
		\end{alertblock}
	\end{frame}
	
	\begin{frame}{Q262: Alert Fatigue Risk}
		\begin{block}{Question}
			What is the MAIN risk of alert fatigue?
		\end{block}
		\begin{enumerate}[(A)]
			\item Fewer alerts
			\item Critical alerts missed among noise
			\item Lower cost
			\item Higher uptime
		\end{enumerate}
	\end{frame}
	
	\begin{frame}{Q262: Explanation}
		\begin{block}{Correct Answer: B}
			\textbf{Critical alerts missed:}
			\begin{itemize}
				\item Analysts ignore alerts
				\item Detection misses real incidents
				\item Tune and prioritize
			\end{itemize}
		\end{block}
		\begin{alertblock}{Critical Trap}
			Alert fatigue causes real incidents to slip.
		\end{alertblock}
	\end{frame}
	
	\begin{frame}{Q263: Threat Hunting Benefit}
		\begin{block}{Question}
			Which benefit does proactive threat hunting provide?
		\end{block}
		\begin{enumerate}[(A)]
			\item Cheaper licensing
			\item Detection of threats that evade automated controls
			\item Fewer analysts needed
			\item Fewer logs
		\end{enumerate}
	\end{frame}
	
	\begin{frame}{Q263: Explanation}
		\begin{block}{Correct Answer: B}
			\textbf{Hunting finds hidden threats:}
			\begin{itemize}
				\item Finds hidden threats
				\item Human-led investigation
				\item Complements automation
			\end{itemize}
		\end{block}
		\begin{alertblock}{Critical Trap}
			Hunting complements but does not replace automation.
		\end{alertblock}
	\end{frame}
	
	\begin{frame}{Q264: Threat Intelligence Feeds}
		\begin{block}{Question}
			How should threat intelligence feeds be used?
		\end{block}
		\begin{enumerate}[(A)]
			\item Blindly block all indicators
			\item Contextualized use with validation for relevance
			\item Ignore all feeds
			\item Store for marketing
		\end{enumerate}
	\end{frame}
	
	\begin{frame}{Q264: Explanation}
		\begin{block}{Correct Answer: B}
			\textbf{Use feeds with context:}
			\begin{itemize}
				\item Not all feeds are accurate
				\item Context is required
				\item Blocking without context harms business
			\end{itemize}
		\end{block}
		\begin{alertblock}{Critical Trap}
			Not every threat indicator applies to your environment.
		\end{alertblock}
	\end{frame}
	
	\begin{frame}{Q265: IR Plan Testing}
		\begin{block}{Question}
			Why should incident response plans be tested?
		\end{block}
		\begin{enumerate}[(A)]
			\item Compliance only
			\item To verify roles, communications, and procedures work in practice
			\item To slow downtime
			\item To add cost
		\end{enumerate}
	\end{frame}
	
	\begin{frame}{Q265: Explanation}
		\begin{block}{Correct Answer: B}
			\textbf{Testing reveals gaps:}
			\begin{itemize}
				\item Tests reveal gaps
				\item Roles must be clear
				\item Communication tested
			\end{itemize}
		\end{block}
		\begin{alertblock}{Critical Trap}
			Untested IR plans fail in real incidents.
		\end{alertblock}
	\end{frame}
	
	\begin{frame}{Q266: Tabletop Value}
		\begin{block}{Question}
			What is the main value of a tabletop IR exercise?
		\end{block}
		\begin{enumerate}[(A)]
			\item Cheaper
			\item Tests decision-making, roles, and communication without production impact
			\item Faster recovery
			\item No risk
		\end{enumerate}
	\end{frame}
	
	\begin{frame}{Q266: Explanation}
		\begin{block}{Correct Answer: B}
			\textbf{Tabletop tests decisions:}
			\begin{itemize}
				\item Decision testing
				\item Role clarity
				\item Communication
			\end{itemize}
		\end{block}
		\begin{alertblock}{Critical Trap}
			Tabletop validates decision-making, not technology only.
		\end{alertblock}
	\end{frame}
	
	\begin{frame}{Q267: Incident Lessons}
		\begin{block}{Question}
			Why should lessons be captured after incidents?
		\end{block}
		\begin{enumerate}[(A)]
			\item Blame
			\item Improve processes, controls, and response capability
			\item Close tickets faster
			\item Marketing
		\end{enumerate}
	\end{frame}
	
	\begin{frame}{Q267: Explanation}
		\begin{block}{Correct Answer: B}
			\textbf{Improve for the future:}
			\begin{itemize}
				\item Continuous improvement
				\item Prevent recurrence
				\item Improve response
			\end{itemize}
		\end{block}
		\begin{alertblock}{Critical Trap}
			Lessons without action are wasted learning.
		\end{alertblock}
	\end{frame}
	
	\begin{frame}{Q268: Forensics Preservation}
		\begin{block}{Question}
			Why is forensic preservation important during an incident?
		\end{block}
		\begin{enumerate}[(A)]
			\item Marketing
			\item Supports investigation, legal action, and root cause analysis
			\item Faster recovery
			\item Cheaper
		\end{enumerate}
	\end{frame}
	
	\begin{frame}{Q268: Explanation}
		\begin{block}{Correct Answer: B}
			\textbf{Preserve evidence:}
			\begin{itemize}
				\item Preserve evidence
				\item Chain of custody
				\item Supports regulators and legal
			\end{itemize}
		\end{block}
		\begin{alertblock}{Critical Trap}
			Poor preservation destroys evidence.
		\end{alertblock}
	\end{frame}
	
	\begin{frame}{Q269: Vendor Breach Communication}
		\begin{block}{Question}
			What should happen when a vendor notifies of a breach?
		\end{block}
		\begin{enumerate}[(A)]
			\item Ignore
			\item Verify scope, assess impact, and coordinate response and notification
			\item Blame the vendor
			\item Wait for public disclosure
		\end{enumerate}
	\end{frame}
	
	\begin{frame}{Q269: Explanation}
		\begin{block}{Correct Answer: B}
			\textbf{Assess and coordinate:}
			\begin{itemize}
				\item Scope assessment
				\item Impact analysis
				\item Regulatory coordination
			\end{itemize}
		\end{block}
		\begin{alertblock}{Critical Trap}
			Vendor breaches are your breaches to manage.
		\end{alertblock}
	\end{frame}
	
	\begin{frame}{Q270: Patch Deployment Risk}
		\begin{block}{Question}
			What is a key risk of rapid patch deployment without testing?
		\end{block}
		\begin{enumerate}[(A)]
			\item Cheaper
			\item Patching can break production systems
			\item Faster
			\item Lower risk
		\end{enumerate}
	\end{frame}
	
	\begin{frame}{Q270: Explanation}
		\begin{block}{Correct Answer: B}
			\textbf{Untested patches break production:}
			\begin{itemize}
				\item Patches can cause regressions
				\item Staged rollout reduces risk
				\item Emergency bypass for zero-day threats
			\end{itemize}
		\end{block}
		\begin{alertblock}{Critical Trap}
			Patching without testing trades one risk for another.
		\end{alertblock}
	\end{frame}
	
	\begin{frame}{Q271: Secure Configuration Baseline}
		\begin{block}{Question}
			What is a secure configuration baseline?
		\end{block}
		\begin{enumerate}[(A)]
			\item Vendor default
			\item Documented secure settings derived from standards and validated
			\item Marketing
			\item Random settings
		\end{enumerate}
	\end{frame}
	
	\begin{frame}{Q271: Explanation}
		\begin{block}{Correct Answer: B}
			\textbf{Baseline is documented secure state:}
			\begin{itemize}
				\item Documented secure state
				\item Based on CIS, NIST, or vendor
				\item Drift detected against baseline
			\end{itemize}
		\end{block}
		\begin{alertblock}{Critical Trap}
			Defaults are rarely secure.
		\end{alertblock}
	\end{frame}
	
	\begin{frame}{Q272: Data Backup vs Archive}
		\begin{block}{Question}
			What is the difference between backup and archive?
		\end{block}
		\begin{enumerate}[(A)]
			\item They are the same
			\item Backup is for recovery; archive is for retention and reference
			\item Archive is faster
			\item Backup is cheaper
		\end{enumerate}
	\end{frame}
	
	\begin{frame}{Q272: Explanation}
		\begin{block}{Correct Answer: B}
			\textbf{Backup vs archive:}
			\begin{itemize}
				\item Backup for restore
				\item Archive for long-term retention
				\item Different retention policies
			\end{itemize}
		\end{block}
		\begin{alertblock}{Critical Trap}
			Backup is not archive; archive is not backup.
		\end{alertblock}
	\end{frame}
	
	\begin{frame}{Q273: Immutable Logs}
		\begin{block}{Question}
			Why should logs be immutable?
		\end{block}
		\begin{enumerate}[(A)]
			\item Cheaper
			\item Prevent tampering and preserve evidence
			\item Faster search
			\item Less storage
		\end{enumerate}
	\end{frame}
	
	\begin{frame}{Q273: Explanation}
		\begin{block}{Correct Answer: B}
			\textbf{Immutable logs preserve evidence:}
			\begin{itemize}
				\item Tamper-evident
				\item Supports forensics and audit
				\item Write-once storage
			\end{itemize}
		\end{block}
		\begin{alertblock}{Critical Trap}
			Logs that can be modified cannot be trusted.
		\end{alertblock}
	\end{frame}
	
	\begin{frame}{Q274: Encryption Key Custody}
		\begin{block}{Question}
			What is the biggest risk of poor key custody?
		\end{block}
		\begin{enumerate}[(A)]
			\item Cheaper
			\item Encryption becomes ineffective if keys are compromised
			\item Slower
			\item Smaller files
		\end{enumerate}
	\end{frame}
	
	\begin{frame}{Q274: Explanation}
		\begin{block}{Correct Answer: B}
			\textbf{Compromised keys break encryption:}
			\begin{itemize}
				\item Keys protect encryption
				\item Poor custody breaks security
				\item Key management critical
			\end{itemize}
		\end{block}
		\begin{alertblock}{Critical Trap}
			Encryption without key custody is fake security.
		\end{alertblock}
	\end{frame}
	
	\begin{frame}{Q275: Access Segregation Risk}
		\begin{block}{Question}
			What risk does an effective segregation of duties prevent?
		\end{block}
		\begin{enumerate}[(A)]
			\item Slower work
			\item Fraud and errors from concentrated control
			\item Cheaper
			\item Faster
		\end{enumerate}
	\end{frame}
	
	\begin{frame}{Q275: Explanation}
		\begin{block}{Correct Answer: B}
			\textbf{SoD prevents concentrated control:}
			\begin{itemize}
				\item Prevent concentration
				\item Reduce fraud
				\item Detect errors
			\end{itemize}
		\end{block}
		\begin{alertblock}{Critical Trap}
			SoD prevents abuse of concentrated authority.
		\end{alertblock}
	\end{frame}
	
	\begin{frame}{Q276: SIEM False Positives}
		\begin{block}{Question}
			What is the best way to reduce SIEM false positives?
		\end{block}
		\begin{enumerate}[(A)]
			\item Disable all rules
			\item Tune correlation rules with business context
			\item Ignore all alerts
			\item Reduce data
		\end{enumerate}
	\end{frame}
	
	\begin{frame}{Q276: Explanation}
		\begin{block}{Correct Answer: B}
			\textbf{Tune rules:}
			\begin{itemize}
				\item Tune rules
				\item Baseline behavior
				\item Add context
			\end{itemize}
		\end{block}
		\begin{alertblock}{Critical Trap}
			Disabling rules removes visibility.
		\end{alertblock}
	\end{frame}
	
	\begin{frame}{Q277: Access Token Risk}
		\begin{block}{Question}
			What risk do stolen access tokens create?
		\end{block}
		\begin{enumerate}[(A)]
			\item Cheaper
			\item Attackers can impersonate users without credentials
			\item Slower
			\item Smaller
		\end{enumerate}
	\end{frame}
	
	\begin{frame}{Q277: Explanation}
		\begin{block}{Correct Answer: B}
			\textbf{Tokens act as credentials:}
			\begin{itemize}
				\item Tokens act as credentials
				\item Short lifetime reduces damage
				\item Monitor and revoke
			\end{itemize}
		\end{block}
		\begin{alertblock}{Critical Trap}
			Tokens must be protected like passwords.
		\end{alertblock}
	\end{frame}
	
	\begin{frame}{Q278: Passwordless Authentication}
		\begin{block}{Question}
			What is the main benefit of passwordless authentication?
		\end{block}
		\begin{enumerate}[(A)]
			\item Cheaper
			\item Reduces phishing and credential compromise risk
			\item Slower
			\item Fewer users
		\end{enumerate}
	\end{frame}
	
	\begin{frame}{Q278: Explanation}
		\begin{block}{Correct Answer: B}
			\textbf{Passwordless reduces phishing:}
			\begin{itemize}
				\item No reusable password
				\item FIDO2 keys, biometrics
				\item Phishing resistant
			\end{itemize}
		\end{block}
		\begin{alertblock}{Critical Trap}
			Passwordless reduces phishing, not all attacks.
		\end{alertblock}
	\end{frame}
	
	\begin{frame}{Q279: Privacy by Design Trade-off}
		\begin{block}{Question}
			What trade-off exists with privacy by design?
		\end{block}
		\begin{enumerate}[(A)]
			\item No trade-off
			\item Some features may become harder to deliver due to minimized data collection
			\item Higher quality
			\item Faster
		\end{enumerate}
	\end{frame}
	
	\begin{frame}{Q279: Explanation}
		\begin{block}{Correct Answer: B}
			\textbf{Limited data constrains features:}
			\begin{itemize}
				\item Limited data constrains features
				\item Requires design work
				\item Long-term trust benefit
			\end{itemize}
		\end{block}
		\begin{alertblock}{Critical Trap}
			Privacy by design trades convenience for trust.
		\end{alertblock}
	\end{frame}
	
	\begin{frame}{Q280: Data Subject Access Request Timelines}
		\begin{block}{Question}
			What is the typical GDPR time frame for responding to a data subject access request?
		\end{block}
		\begin{enumerate}[(A)]
			\item 24 hours
			\item 30 days, extendable in some cases
			\item 1 year
			\item 5 years
		\end{enumerate}
	\end{frame}
	
	\begin{frame}{Q280: Explanation}
		\begin{block}{Correct Answer: B}
			\textbf{One month response:}
			\begin{itemize}
				\item One month
				\item Extendable for complexity
				\item Verify identity first
			\end{itemize}
		\end{block}
		\begin{alertblock}{Critical Trap}
			GDPR DSAR response is one month, not one week.
		\end{alertblock}
	\end{frame}
	
	\begin{frame}{Q281: Data Subject Rights Refusal}
		\begin{block}{Question}
			When may an organization refuse a data subject request?
		\end{block}
		\begin{enumerate}[(A)]
			\item Never
			\item Legal exemptions, third-party rights, or manifestly unfounded or excessive requests
			\item Whenever convenient
			\item Because it costs money
		\end{enumerate}
	\end{frame}
	
	\begin{frame}{Q281: Explanation}
		\begin{block}{Correct Answer: B}
			\textbf{Legal justification required:}
			\begin{itemize}
				\item Limited legal exemptions
				\item Document rationale
				\item Must explain refusal
			\end{itemize}
		\end{block}
		\begin{alertblock}{Critical Trap}
			Refusals require legal justification.
		\end{alertblock}
	\end{frame}
	
	\begin{frame}{Q282: HIPAA Breach Risk Assessment}
		\begin{block}{Question}
			What must a HIPAA covered entity do after discovering a breach?
		\end{block}
		\begin{enumerate}[(A)]
			\item Ignore
			\item Conduct a risk assessment to determine notification requirement
			\item Notify all individuals
			\item Publish a press release
		\end{enumerate}
	\end{frame}
	
	\begin{frame}{Q282: Explanation}
		\begin{block}{Correct Answer: B}
			\textbf{Documented risk assessment:}
			\begin{itemize}
				\item Risk assessment
				\item Notification requirements depend on results
				\item Presumption of breach unless low probability of compromise
			\end{itemize}
		\end{block}
		\begin{alertblock}{Critical Trap}
			HIPAA requires a documented risk assessment.
		\end{alertblock}
	\end{frame}
	
	\begin{frame}{Q283: CCPA Opt-Out}
		\begin{block}{Question}
			What does CCPA require regarding opt-out of the sale of personal information?
		\end{block}
		\begin{enumerate}[(A)]
			\item Nothing
			\item A clear ``Do Not Sell My Personal Information'' mechanism
			\item Only verbal requests
			\item Only via mail
		\end{enumerate}
	\end{frame}
	
	\begin{frame}{Q283: Explanation}
		\begin{block}{Correct Answer: B}
			\textbf{CCPA requires opt-out mechanism:}
			\begin{itemize}
				\item Opt-out right
				\item Clear mechanism
				\item Enforced
			\end{itemize}
		\end{block}
		\begin{alertblock}{Critical Trap}
			CCPA requires opt-out mechanisms for sale.
		\end{alertblock}
	\end{frame}
	
	\begin{frame}{Q284: Encryption Strength}
		\begin{block}{Question}
			What determines the strength of encryption?
		\end{block}
		\begin{enumerate}[(A)]
			\item Amount of data
			\item Algorithm and key length combined
			\item Storage location
			\item User count
		\end{enumerate}
	\end{frame}
	
	\begin{frame}{Q284: Explanation}
		\begin{block}{Correct Answer: B}
			\textbf{Algorithm plus key length:}
			\begin{itemize}
				\item Strong algorithm
				\item Sufficient key length
				\item Key management
			\end{itemize}
		\end{block}
		\begin{alertblock}{Critical Trap}
			Encryption strength needs both algorithm and key length.
		\end{alertblock}
	\end{frame}
	
	\begin{frame}{Q285: Digital Signature Verification}
		\begin{block}{Question}
			How does a recipient verify a digital signature?
		\end{block}
		\begin{enumerate}[(A)]
			\item Using the sender's private key
			\item Using the sender's public key
			\item Using the recipient's private key
			\item Using the recipient's public key
		\end{enumerate}
	\end{frame}
	
	\begin{frame}{Q285: Explanation}
		\begin{block}{Correct Answer: B}
			\textbf{Verify with sender public key:}
			\begin{itemize}
				\item Sender signs with private key
				\item Recipient verifies with public key
				\item Ensures authenticity
			\end{itemize}
		\end{block}
		\begin{alertblock}{Critical Trap}
			Verify signatures with the sender's public key.
		\end{alertblock}
	\end{frame}
	
	\begin{frame}{Q286: Security Awareness for Executives}
		\begin{block}{Question}
			Why do executives need specialized security awareness?
		\end{block}
		\begin{enumerate}[(A)]
			\item They don't need training
			\item They are prime targets for spear phishing and social engineering
			\item They already know everything
			\item It is optional
		\end{enumerate}
	\end{frame}
	
	\begin{frame}{Q286: Explanation}
		\begin{block}{Correct Answer: B}
			\textbf{Executives are top targets:}
			\begin{itemize}
				\item High-value targets
				\item Specialized risks
				\item Frequent attacks
			\end{itemize}
		\end{block}
		\begin{alertblock}{Critical Trap}
			Executives are top phishing targets.
		\end{alertblock}
	\end{frame}
	
	\begin{frame}{Q287: Whaling Attack}
		\begin{block}{Question}
			What is a whaling attack?
		\end{block}
		\begin{enumerate}[(A)]
			\item Attack on whales
			\item Targeted phishing on senior executives
			\item Attack on routers
			\item Network attack
		\end{enumerate}
	\end{frame}
	
	\begin{frame}{Q287: Explanation}
		\begin{block}{Correct Answer: B}
			\textbf{Whaling targets executives:}
			\begin{itemize}
				\item High-value target
				\item Spear phishing variant
				\item Executive protection matters
			\end{itemize}
		\end{block}
		\begin{alertblock}{Critical Trap}
			Whaling = executives targeted by phishing.
		\end{alertblock}
	\end{frame}
	
	\begin{frame}{Q288: Tabletop for Board}
		\begin{block}{Question}
			Why would a board participate in a cyber tabletop exercise?
		\end{block}
		\begin{enumerate}[(A)]
			\item Entertainment
			\item Test crisis decision-making and board oversight during a major incident
			\item Compliance only
			\item Marketing
		\end{enumerate}
	\end{frame}
	
	\begin{frame}{Q288: Explanation}
		\begin{block}{Correct Answer: B}
			\textbf{Board tests crisis oversight:}
			\begin{itemize}
				\item Governance during crisis
				\item Decision rights
				\item Board roles
			\end{itemize}
		\end{block}
		\begin{alertblock}{Critical Trap}
			Boards need to practice oversight under stress.
		\end{alertblock}
	\end{frame}
	
	\begin{frame}{Q289: Risk of Unpatched Middleware}
		\begin{block}{Question}
			Why is unpatched middleware particularly dangerous?
		\end{block}
		\begin{enumerate}[(A)]
			\item Marketing
			\item Middleware sits between applications and databases and can expose high-value systems
			\item Cheaper
			\item Faster
		\end{enumerate}
	\end{frame}
	
	\begin{frame}{Q289: Explanation}
		\begin{block}{Correct Answer: B}
			\textbf{Middleware has amplified impact:}
			\begin{itemize}
				\item High-value position
				\item Broad impact
				\item Critical patch priority
			\end{itemize}
		\end{block}
		\begin{alertblock}{Critical Trap}
			Middleware vulnerabilities have amplified impact.
		\end{alertblock}
	\end{frame}
	
	\begin{frame}{Q290: API Key Storage Risk}
		\begin{block}{Question}
			What is the main risk of API keys hardcoded in code?
		\end{block}
		\begin{enumerate}[(A)]
			\item Cheaper
			\item Leakage through code repositories or logs
			\item Faster
			\item Lower risk
		\end{enumerate}
	\end{frame}
	
	\begin{frame}{Q290: Explanation}
		\begin{block}{Correct Answer: B}
			\textbf{Code leaks secrets:}
			\begin{itemize}
				\item Code is copied
				\item Repos leak
				\item Use secret managers
			\end{itemize}
		\end{block}
		\begin{alertblock}{Critical Trap}
			Hardcoded secrets are exposed to anyone with code access.
		\end{alertblock}
	\end{frame}
	
	\begin{frame}{Q291: Continuous Compliance}
		\begin{block}{Question}
			What is continuous compliance in cloud environments?
		\end{block}
		\begin{enumerate}[(A)]
			\item Annual audits
			\item Continuous automated monitoring against compliance rules
			\item Quarterly reviews
			\item Never
		\end{enumerate}
	\end{frame}
	
	\begin{frame}{Q291: Explanation}
		\begin{block}{Correct Answer: B}
			\textbf{Continuous checks:}
			\begin{itemize}
				\item Automated checks
				\item Real-time detection
				\item Remediation
			\end{itemize}
		\end{block}
		\begin{alertblock}{Critical Trap}
			Continuous compliance detects drift continuously.
		\end{alertblock}
	\end{frame}
	
	\begin{frame}{Q292: CAIQ vs SOC 2}
		\begin{block}{Question}
			What is the difference between CAIQ and SOC 2?
		\end{block}
		\begin{enumerate}[(A)]
			\item They are the same
			\item CAIQ is a self-assessment questionnaire; SOC 2 is independent audit
			\item CAIQ is an audit
			\item SOC 2 is a questionnaire
		\end{enumerate}
	\end{frame}
	
	\begin{frame}{Q292: Explanation}
		\begin{block}{Correct Answer: B}
			\textbf{Self-report vs audit:}
			\begin{itemize}
				\item CAIQ is self-report
				\item SOC 2 is audited
				\item Different assurance levels
			\end{itemize}
		\end{block}
		\begin{alertblock}{Critical Trap}
			CAIQ is self-assessment; SOC 2 is independent.
		\end{alertblock}
	\end{frame}
	
	\begin{frame}{Q293: SaaS Misconfiguration}
		\begin{block}{Question}
			What is the most common cause of SaaS data exposure?
		\end{block}
		\begin{enumerate}[(A)]
			\item Vendor vulnerability
			\item Customer misconfiguration of permissions or sharing
			\item Cheaper
			\item Storage limits
		\end{enumerate}
	\end{frame}
	
	\begin{frame}{Q293: Explanation}
		\begin{block}{Correct Answer: B}
			\textbf{Customer misconfigurations:}
			\begin{itemize}
				\item Permissions
				\item Public sharing
				\item Shared responsibility
			\end{itemize}
		\end{block}
		\begin{alertblock}{Critical Trap}
			Most SaaS breaches are customer misconfigurations.
		\end{alertblock}
	\end{frame}
	
	\begin{frame}{Q294: Identity in Cloud}
		\begin{block}{Question}
			What is the most critical control for cloud security?
		\end{block}
		\begin{enumerate}[(A)]
			\item Firewalls
			\item Strong identity and access management
			\item Antivirus
			\item Backups
		\end{enumerate}
	\end{frame}
	
	\begin{frame}{Q294: Explanation}
		\begin{block}{Correct Answer: B}
			\textbf{Cloud security starts with IAM:}
			\begin{itemize}
				\item IAM is fundamental
				\item Least privilege
				\item MFA
			\end{itemize}
		\end{block}
		\begin{alertblock}{Critical Trap}
			Cloud security starts with IAM.
		\end{alertblock}
	\end{frame}
	
	\begin{frame}{Q295: Cloud Encryption}
		\begin{block}{Question}
			What is the main challenge with cloud encryption?
		\end{block}
		\begin{enumerate}[(A)]
			\item Slower
			\item Key management and provider access
			\item Cheaper
			\item Less storage
		\end{enumerate}
	\end{frame}
	
	\begin{frame}{Q295: Explanation}
		\begin{block}{Correct Answer: B}
			\textbf{Key custody matters:}
			\begin{itemize}
				\item Key custody
				\item Provider access
				\item BYOK and HYOK options
			\end{itemize}
		\end{block}
		\begin{alertblock}{Critical Trap}
			Cloud encryption depends on key control.
		\end{alertblock}
	\end{frame}
	
	\begin{frame}{Q296: Data Loss in Cloud}
		\begin{block}{Question}
			What causes data loss in cloud environments?
		\end{block}
		\begin{enumerate}[(A)]
			\item Provider failures only
			\item Misconfiguration, deletion by attacker, insufficient backups
			\item Cheaper storage
			\item Faster networks
		\end{enumerate}
	\end{frame}
	
	\begin{frame}{Q296: Explanation}
		\begin{block}{Correct Answer: B}
			\textbf{Often customer-caused:}
			\begin{itemize}
				\item Misconfigurations
				\item Accidental deletion
				\item Backups matter
			\end{itemize}
		\end{block}
		\begin{alertblock}{Critical Trap}
			Cloud data loss is often customer-caused.
		\end{alertblock}
	\end{frame}
	
	\begin{frame}{Q297: Security Metrics}
		\begin{block}{Question}
			Which security metric is most valuable to executives?
		\end{block}
		\begin{enumerate}[(A)]
			\item Number of CVEs fixed
			\item Trend in risk exposure against appetite and business impact
			\item Firewall rule count
			\item Number of employees
		\end{enumerate}
	\end{frame}
	
	\begin{frame}{Q297: Explanation}
		\begin{block}{Correct Answer: B}
			\textbf{Business context matters:}
			\begin{itemize}
				\item Business impact matters
				\item Appetite alignment
				\item Trends show direction
			\end{itemize}
		\end{block}
		\begin{alertblock}{Critical Trap}
			Executives need business context, not raw counts.
		\end{alertblock}
	\end{frame}
	
	\begin{frame}{Q298: Risk Acceptance in Security}
		\begin{block}{Question}
			Who should accept residual security risk?
		\end{block}
		\begin{enumerate}[(A)]
			\item The security team
			\item The business owner, informed by security
			\item The vendor
			\item The auditor
		\end{enumerate}
	\end{frame}
	
	\begin{frame}{Q298: Explanation}
		\begin{block}{Correct Answer: B}
			\textbf{Business owns risk:}
			\begin{itemize}
				\item Business owns risk
				\item Security advises
				\item Documented acceptance
			\end{itemize}
		\end{block}
		\begin{alertblock}{Critical Trap}
			Security advises; business owners accept risk.
		\end{alertblock}
	\end{frame}
	
	\begin{frame}{Q299: Security Program Maturity}
		\begin{block}{Question}
			What BEST indicates high security program maturity?
		\end{block}
		\begin{enumerate}[(A)]
			\item More tools
			\item Risk-based decisions, metrics, and continuous improvement
			\item Larger budgets
			\item Fewer incidents
		\end{enumerate}
	\end{frame}
	
	\begin{frame}{Q299: Explanation}
		\begin{block}{Correct Answer: B}
			\textbf{Maturity is outcomes:}
			\begin{itemize}
				\item Risk-based
				\item Measured
				\item Improving
			\end{itemize}
		\end{block}
		\begin{alertblock}{Critical Trap}
			Maturity is about outcomes and improvement, not tools.
		\end{alertblock}
	\end{frame}
	
	\begin{frame}{Q300: Ultimate Security Truth}
		\begin{block}{Question}
			What is the ultimate truth about security?
		\end{block}
		\begin{enumerate}[(A)]
			\item Security can be perfect
			\item Security can never be perfect; it is continuous risk management
			\item Security is a one-time project
			\item Security is a purchase
		\end{enumerate}
	\end{frame}
	
	\begin{frame}{Q300: Explanation}
		\begin{block}{Correct Answer: B}
			\textbf{Security is continuous:}
			\begin{itemize}
				\item No perfect security
				\item Continuous risk management
				\item Balance risk and cost
			\end{itemize}
		\end{block}
		\begin{alertblock}{Critical Trap}
			Security is continuous, not a one-time achievement.
		\end{alertblock}
	\end{frame}
	
	\begin{frame}{Exam Mindset Summary}
		\begin{block}{Remember}
			\begin{itemize}
				\item \textbf{CIA triad:} Confidentiality, Integrity, Availability
				\item \textbf{Signatures} provide integrity, authenticity, non-repudiation
				\item \textbf{Encryption is reversible; hashing is one-way}
				\item \textbf{Privacy} is about rights; \textbf{security} is about protection
				\item \textbf{Zero Trust:} verify every access
				\item \textbf{DLP, PAM, SIEM, WAF, HSM, KMS:} know what each does
				\item \textbf{Cloud:} shared responsibility; data location matters
				\item \textbf{Maturity:} continuous improvement and metrics
				\item \textbf{Security is continuous risk management}
			\end{itemize}
		\end{block}
	\end{frame}
	
	\begin{frame}{Thank You}
		\begin{center}
			{\Huge \textbf{Good Luck on Your CRISC Exam!}}
		\end{center}
	\end{frame}
	
\end{document}